\documentclass[11pt]{article}
\usepackage{amsmath,amssymb,amsthm,booktabs,geometry,graphicx}
\usepackage[utf8]{inputenc}
\usepackage[T1]{fontenc}
\usepackage[french]{babel}
\usepackage[round,authoryear]{natbib}
\usepackage[colorlinks=true,linkcolor=black,citecolor=black,urlcolor=blue!50!black]{hyperref}
\geometry{margin=1in}
\setlength{\emergencystretch}{2em}

\newtheorem{proposition}{Proposition}
\newtheorem{corollary}{Corollaire}
\newtheorem{assumption}{Hypothèse}
\theoremstyle{definition}
\newtheorem{postulate}{Postulat}
\newtheorem{definition}{Définition}
\theoremstyle{remark}
\newtheorem*{remark}{Remarque}
\renewcommand{\proofname}{Démonstration}

\newcommand{\JJ}{\mathcal{J}}
\newcommand{\DD}{\mathcal{D}}
\newcommand{\Dun}{\mathcal{D}_1}
\newcommand{\KL}[2]{D_{KL}\!\left(#1\,\Vert\,#2\right)}

\title{Compléments théoriques\\[4pt]
\large Fondements informationnels d'un indicateur entropique\\de stress systémique : ce que le papier principal laisse de côté}
\author{Annexe au papier \og Sauts d'entropie et relaxation\\de l'information sur les marchés financiers \fg{}}
\date{Octobre 2026}

\begin{document}
\maketitle

\begin{abstract}
\noindent Le papier principal est recentré sur un seul objet destiné aux
banques centrales : le canal de dépendance $-\tfrac12\log\det R$ d'un panel
multi-marchés, sa décomposition exacte entre blocs et à l'intérieur des blocs,
et l'audit en nats des scénarios de stress publiés. Ce document rassemble tout
ce que ce recentrage retire du corps du texte, pour servir d'annexe ou être
fourni sur demande. On y trouve le pont qui mène du maximum d'entropie en
temps-volume au mouvement brownien géométrique, la décomposition de l'entropie
en trois canaux et l'indice structurel $\JJ=\KL{p}{\prod_i q_i}$, la dynamique
de saut et diffusion (un théorème, deux postulats, la structure
d'Ornstein--Uhlenbeck imposée par le théorème de Doob), le programme empirique
H1--H8 avec son verdict face aux témoins calibrés, et la géométrie complète des
scénarios en boule de Kullback--Leibler. Deux résultats sont nouveaux par
rapport à la version arXiv~v1. Le premier est une identité de décomposition par
blocs du canal de dépendance, qui relie ce document au papier principal. Le
second montre que le relèvement d'un scénario de portefeuille à la dimension
pleine qui minimise la divergence ambiante coûte exactement le rayon $\eta$, et
qu'il s'interprète actif par actif (choc porté par la covariance de chaque actif
avec le portefeuille), alors que la dilatation uniforme par défaut coûte $N$~fois
plus en variance. Ce second point corrige une remarque de la v1.
\end{abstract}

\tableofcontents

\section*{Avant-propos : ce que contient ce document, et pourquoi}
\addcontentsline{toc}{section}{Avant-propos}

\paragraph{L'articulation avec le papier principal.} L'évaluation d'octobre 2026
a conclu qu'un seul résultat de la version arXiv~v1 franchit les deux témoins
calibrés sur les deux panels : l'écart de régime du canal de dépendance
$H^{dep}=\tfrac12\log\det R$. La phase~A l'a confirmé sous trois formes (brute,
dévolatilisée à la Forbes--Rigobon, en rangs) et a montré que le ratio
d'absorption de \citet{kritzman2011}, mesure de référence de la même famille,
échoue au témoin par blocs. Le papier principal se reconstruit donc autour d'un
\og CISS entropique \fg{} au sens de \citet{hollo2012} : canal de dépendance d'un
panel multi-marchés, décomposition intra/inter-blocs, validation à la manière
des banques centrales, et audit des scénarios publiés en nats.

Tout le reste sort du corps du texte. Ce n'est pas parce que c'est faux : les
résultats mathématiques restent exacts. C'est parce que ces éléments ne sont
pas nécessaires à l'indicateur, ou parce qu'ils ne survivent pas aux témoins et
ne peuvent donc pas fonder un outil de surveillance. Les rassembler ici a deux
usages. Un rapporteur qui demande d'où vient le canal de dépendance, et pourquoi
c'est lui et pas l'indice structurel complet, trouve la réponse. Un lecteur qui
veut reprendre la théorie de la dynamique de l'information trouve la chaîne
entière, avec le statut exact de chaque maillon.

\paragraph{Statut de chaque énoncé.} Le document distingue trois catégories et
les signale au fil du texte : les \emph{théorèmes} (exacts sous leurs
hypothèses), les \emph{postulats} (hypothèses économiques substantielles,
testables), et les \emph{constats empiriques} (avec le témoin qu'ils
franchissent ou non). Le tableau~\ref{tab:statut} en donne la carte.

\begin{table}[htbp]
\centering
\small
\begin{tabular}{@{}p{5.0cm}p{3.0cm}p{4.6cm}l@{}}
\toprule
Élément & Statut & Verdict empirique & Section \\
\midrule
Maximum d'entropie gaussien & théorème & --- & \ref{sec:bridge} \\
Pont volume (MDH, brownien subordonné) & argument sous hypothèses & non testé ici & \ref{sec:bridge} \\
Décomposition en trois canaux & identité exacte & --- & \ref{sec:channels} \\
Indice structurel $\JJ$ comme divergence & identité exacte & étendue franchit les deux témoins & \ref{sec:channels} \\
Sauts d'entropie vers le bas & théorème (en moyenne) & non estimable directement & \ref{sec:dynamics} \\
Relaxation (P1) & postulat & demi-vie non identifiée & \ref{sec:dynamics} \\
Contrainte unilatérale (P2) & postulat & non établie (H5) & \ref{sec:dynamics} \\
Couplage queues--corrélation & conséquence & non identifié (H6) & \ref{sec:dynamics} \\
Canal de dépendance sous stress & constat & \textbf{franchit les deux témoins} & \ref{sec:program}, \ref{sec:blocks} \\
Boule KL, $1$~nat $=\sqrt2\,\sigma$ & théorème & --- & \ref{sec:stress} \\
Relèvement minimal (rang un) & théorème (nouveau) & --- & \ref{sec:stress} \\
Calibration du rayon sur le flux & règle & trop douce face à l'histoire & \ref{sec:stress} \\
Décomposition par blocs & identité (nouvelle ici) & objet du papier principal & \ref{sec:blocks} \\
\bottomrule
\end{tabular}
\caption{Carte des énoncés du document et de leur statut.}
\label{tab:statut}
\end{table}

\paragraph{Origine des textes.} Les sections~\ref{sec:bridge} à
\ref{sec:dynamics} et les annexes~\ref{app:units} à \ref{app:mi} reprennent, avec
des retouches de renvois, le texte de la version française de référence
(\texttt{paper/main\_fr.tex}). Les sections~\ref{sec:program} à
\ref{sec:open} sont rédigées pour ce document. Tous les chiffres cités
proviennent du même code (\texttt{python -m src.run\_empirical}, témoins à
deux cents réplications, phase~A) ; aucun n'a été recalculé pour cette annexe.

\section*{Notations}
\addcontentsline{toc}{section}{Notations}

\noindent Deux mises en garde avant le tableau. D'abord, l'entropie
différentielle $h$ \textbf{dépend des unités} : $h(aX)=h(X)+\log|a|$, de sorte
qu'exprimer les rendements en pourcentages plutôt qu'en décimales déplace $h$
d'une constante. Les divergences ($\KL{\cdot}{\cdot}$, $I$, $\DD$, $\JJ$) sont au
contraire invariantes par tout changement de variable inversible : c'est la
raison d'être de l'indice structurel. Ensuite, $K$ désigne un kurtosis et les
lettres \og KL \fg{} sont les initiales de Kullback et Leibler ; il n'y a aucun
lien entre les deux.

\begin{center}
\small
\begin{tabular}{@{}lll@{}}
\toprule
Symbole & Signification & Nature \\
\midrule
$p_t$ & loi conditionnelle vraie des rendements en $t$ & --- \\
$q_t$ & gaussienne de mêmes moyenne et covariance que $p_t$ & --- \\
$h(\cdot)$ & entropie différentielle & nats, \emph{dépend des unités} \\
$\KL{p}{q}$ & divergence de Kullback--Leibler & nats, invariante \\
$\DD_t$ & déficit d'entropie $=\KL{p_t}{q_t}\ge0$ & nats, invariante \\
$\Dun(z)$ & déficit d'un scalaire standardisé (néguentropie) & nats, invariante \\
$H^{vol}_t$ & canal d'échelle $=\sum_i\log\sigma_{i,t}$ & nats, \emph{dépend des unités} \\
$H^{dep}_t$ & canal de dépendance $=\tfrac12\log\det R_t\le0$ & nats, invariante \\
$H^{cov}_t$ & $H^{vol}_t+H^{dep}_t$ ; la constante $\tfrac N2\log2\pi e$ est omise & nats \\
$\JJ_t$ & indice structurel $=\DD_t-H^{dep}_t\ge0$ & nats, invariante \\
$I(X;Y)$ & information mutuelle & nats, invariante \\
$I^{(h)}_t$ & flux d'information $=\KL{p_t}{p_{t-h}}$ & nats \\
$S$, $K$ & asymétrie, kurtosis excédentaire & sans dimension \\
$\Sigma_t=D_tR_tD_t$ & covariance, écarts types, corrélation & --- \\
$\lambda_k$, $N^{\text{eff}}_t$ & valeurs propres de $R_t$, nombre effectif de modes & $1\le N^{\text{eff}}\le N$ \\
$\kappa$, $\sigma_\JJ$ & taux de relaxation et coefficient de diffusion de $\JJ$ & --- \\
$\eta$ & rayon d'un scénario de stress & nats \\
$N$ & nombre d'actifs & --- \\
\bottomrule
\end{tabular}
\end{center}
\section{Du flux d'ordres au prix gaussien : le pont volume}
\label{sec:bridge}

\subsection{Le prix comme flux d'ordres accumulé}

Sur $[t,t+\Delta]$, soit $M$ le nombre de transactions exécutées, la transaction
$k$ portant le volume signé $\epsilon_k v_k$ avec $\epsilon_k\in\{-1,+1\}$ le
sens (à l'initiative de l'acheteur ou du vendeur) et $v_k>0$ la taille. Notons
$V=\sum_k v_k$ le volume total et $Q=\sum_k \epsilon_k v_k$ le flux d'ordres
signé net.

\begin{assumption}[Impact linéaire]\label{as:impact}
Le log-prix répond linéairement au flux d'ordres signé net,
$\Delta X_{t,\Delta} = X_{t+\Delta}-X_t = \lambda\,Q$, avec $\lambda>0$ constant
sur l'intervalle.
\end{assumption}

C'est l'équilibre de \citet{kyle1985} : le teneur de marché, incapable de
distinguer le flux informé du flux non informé, fixe une règle de prix linéaire
et $\lambda$ est l'inverse de la profondeur. C'est la définition du \og prix
comme fonction du volume échangé \fg{} dont a besoin le reste de la section, et
c'est le maillon le plus faible de la chaîne ; la
section~\ref{sec:bridge-caveats} dit ce qui se passe quand il cède.

\begin{assumption}[Signes échangeables et variance proportionnelle au volume]
\label{as:signs}
Conditionnellement à $V$ et à l'information disponible en $t$, la croyance du
marché sur $Q$ n'est contrainte que par ses deux premiers moments, avec
$\mathbb{E}[Q\mid V] = m$ et $\mathrm{Var}(Q\mid V) = \varsigma^2 V$.
\end{assumption}

La forme de la variance est celle qu'on obtient lorsque les signes des
transactions sont échangeables à corrélation bornée et que les tailles ont une
moyenne stable : pour des signes indépendants,
$\mathrm{Var}(Q)=\sum_k v_k^2 \approx \bar v V$. Le contenu économique est que
\emph{le volume, et non le temps, est la mesure naturelle de ce qui peut se
produire}.

\subsection{Maximum d'entropie en temps-volume}

\begin{proposition}[Maximum d'entropie sous contraintes de moments d'ordre un et deux]
\label{prop:maxent}
Parmi les distributions absolument continues sur $\mathbb{R}^N$ de moyenne $\mu$
et de covariance $\Sigma\succ0$, la gaussienne $\mathcal N(\mu,\Sigma)$ maximise
de façon unique l'entropie différentielle, et
$h^{\max}(\Sigma)=\tfrac12\log\left((2\pi e)^N\det\Sigma\right)$.
\end{proposition}

\begin{proof}
Soit $q=\mathcal N(\mu,\Sigma)$ et soit $p$ de mêmes deux premiers moments et
d'entropie différentielle finie. Comme $\log q$ est exactement quadratique, son
espérance ne dépend que de ces moments, donc
$\mathbb{E}_p[\log q]=\mathbb{E}_q[\log q]=-h(q)$. La positivité de l'entropie
relative donne $0\le \KL{p}{q} = -h(p)+h(q)$, d'où $h(p)\le h(q)$, avec égalité
si et seulement si $p=q$ presque partout. L'argument n'exige aucune
différentiabilité de $p$ ni l'existence préalable d'un maximiseur.
\end{proof}

La voie variationnelle --- un lagrangien sur les contraintes de normalisation, de
moyenne et de covariance, dont la condition du premier ordre rend $\log p$
quadratique --- explique \emph{pourquoi} le maximiseur est gaussien, mais c'est
l'argument par la divergence qui établit l'optimalité et l'unicité.

En appliquant la proposition~\ref{prop:maxent} à
l'hypothèse~\ref{as:signs} en dimension un, puis en transportant par
l'hypothèse~\ref{as:impact} :

\begin{equation}
\label{eq:mdh}
\Delta X_{t,\Delta}\ \big|\ V \ \sim\ \mathcal{N}\!\left(\lambda m,\ \lambda^2\varsigma^2 V\right).
\end{equation}

L'équation~\eqref{eq:mdh} est le pont. Trois points méritent d'être soulignés.

D'abord, \emph{la variable de conditionnement est le volume, pas le temps}. Le
maximum d'entropie ne délivre pas une variance ; il délivre la distribution la
moins engageante étant donné une variance. L'hypothèse~\ref{as:signs} fournit
cette variance à partir d'un observable, et cet observable est le volume échangé.
L'équation~\eqref{eq:mdh} est l'hypothèse de mélange de distributions de
\citet{clark1973} et \citet{tauchen1983}, obtenue ici comme un énoncé de maximum
d'entropie plutôt que posée ; \citet{ane2000} montrent empiriquement que les
rendements mesurés en temps-transaction sont proches du gaussien alors que les
rendements en temps calendaire ne le sont pas, ce qui est la prédiction
de~\eqref{eq:mdh}.

Ensuite, \emph{la gaussianité est conditionnelle}. Inconditionnellement,
$\Delta X$ est un mélange d'échelle de lois normales sur le volume aléatoire $V$,
donc leptokurtique même si \eqref{eq:mdh} est exacte. C'est le mécanisme de
subordination de \citet{mandelbrot1967}, et cela signifie que les queues épaisses
ne constituent pas à elles seules une preuve contre le maximum d'entropie.

Enfin, l'extension multivariée est immédiate : avec $N$ actifs,
$\Delta \mathbf{X}\mid V \sim \mathcal N(\boldsymbol\mu V, \Sigma V)$, et la
proposition~\ref{prop:maxent} s'applique telle quelle.

\subsection{Des incréments gaussiens au mouvement brownien puis au MBG}

Soit $\theta_t=\int_0^t v_s\,ds$ l'horloge de volume cumulé.
L'équation~\eqref{eq:mdh} donne la loi d'un incrément ; un processus demande
davantage.

\begin{assumption}[Cohérence temporelle]\label{as:time}
Dans l'horloge de volume, les incréments sur des intervalles disjoints sont
indépendants, les trajectoires sont continues, et la covariance d'un incrément
est proportionnelle au volume écoulé.
\end{assumption}

Sous l'hypothèse~\ref{as:time} à coefficients constants, il existe une matrice
$B$ telle que $BB^{\!\top}=\Sigma$ et $X_\theta = X_0 + \mu\theta + BW_\theta$ en
temps-volume, donc en temps calendaire
\begin{equation}
\label{eq:subordination}
X_t = X_0 + \mu\,\theta_t + B\,W_{\theta_t},
\end{equation}
un mouvement brownien \emph{subordonné} à l'horloge de volume. Si l'intensité de
négociation est constante, $\theta_t=\bar v\,t$ et \eqref{eq:subordination} est
un mouvement brownien ordinaire en temps calendaire ; en exponentiant et en
appliquant Itô,
\begin{equation}
\frac{dS_{i,t}}{S_{i,t}} = \left(a_{i,t}+\tfrac12\sigma_{i,t}^2\right)dt + \sigma_{i,t}\,dW_{i,t},
\end{equation}
de sorte que la dérive du prix n'est pas celle du log-prix --- la même correction
du second ordre qui reparaîtra à la section~\ref{sec:channels} lorsqu'on
différenciera $\log\det R_t$.

Il vaut la peine d'être précis sur ce que fait l'hypothèse~\ref{as:time}. Par le
théorème de \citet{monroe1978}, \emph{toute} martingale locale continue est un
mouvement brownien changé de temps : la subordination seule n'a donc aucun
contenu empirique. Le contenu est dans l'horloge : \eqref{eq:mdh} dit que la
bonne horloge est le volume, et l'hypothèse~\ref{as:time} dit que le marché est
sans mémoire dans cette horloge. C'est le maximum d'entropie qui sélectionne le
mouvement brownien dans la classe qu'autorise le théorème de Monroe.

\subsection{La contraposée}
\label{sec:contrapositive}

La proposition~\ref{prop:maxent} se lit dans les deux sens. Notons $q_t$ la
gaussienne de mêmes moyenne et covariance conditionnelles que $p_t$. Alors

\begin{equation}
\label{eq:deficit}
\DD_t \;\equiv\; h(q_t) - h(p_t) \;=\; \KL{p_t}{q_t} \;\ge\; 0,
\end{equation}
avec égalité si et seulement si $p_t$ est gaussienne. Donc : \emph{les prix sont
non gaussiens exactement dans la mesure où leur distribution conditionnelle n'est
pas d'entropie maximale}, et $\DD_t$ mesure cette mesure. Dans le langage de la
section~\ref{sec:bridge}, un $\DD_t$ positif signifie que des contraintes
au-delà des deux premiers moments sont actives --- que quelqu'un sait quelque
chose que le moment d'ordre deux n'encode pas.

Deux sources produisent un déficit positif, et les séparer importe. Une horloge
de volume aléatoire produit un mélange d'échelle, symétrique et leptokurtique.
Une information directionnelle contraignante produit de l'asymétrie. La
section~\ref{sec:channels} scinde $\DD$ exactement selon cette ligne.

\subsection{Où le pont est faible}
\label{sec:bridge-caveats}

L'hypothèse~\ref{as:impact} est linéaire ; l'impact empirique d'un métaordre est
plus proche d'une racine carrée de la taille, ce qui rend concave l'application
du flux au prix et la distribution de rendements induite moins épaisse en queue
que le flux. L'hypothèse~\ref{as:signs} traite les signes comme échangeables ;
le flux d'ordres est en réalité fortement autocorrélé, ce qui gonfle
$\mathrm{Var}(Q)$ au-dessus de $\varsigma^2 V$ et fait de l'horloge de volume un
substitut imparfait de la variance. Aucune de ces défaillances ne change la
\emph{forme} de l'argument --- le maximum d'entropie sous budget de variance
reste gaussien --- mais toutes deux signifient que l'horloge de volume est une
approximation de l'horloge d'information plutôt que cette horloge même. Nous
utilisons le raisonnement en temps-volume pour motiver le cadre et l'estimation
en temps calendaire dans tout le travail empirique, de sorte qu'aucun résultat
ci-dessous ne dépend de l'exactitude de \eqref{eq:mdh}.

\section{Les trois canaux de l'entropie}
\label{sec:channels}

\subsection{Échelle et dépendance}

Écrivons $\Sigma_t = D_tR_tD_t$ avec
$D_t=\mathrm{diag}(\sigma_{1,t},\dots,\sigma_{N,t})$. Alors
$\det\Sigma_t = \left(\prod_i\sigma_{i,t}^2\right)\det R_t$ et
\begin{equation}
\label{eq:ceiling}
h(q_t) \;=\; \underbrace{\tfrac{N}{2}\log(2\pi e)}_{\text{constante}}
\;+\; \underbrace{\sum_i \log\sigma_{i,t}}_{H^{vol}_t}
\;+\; \underbrace{\tfrac12\log\det R_t}_{H^{dep}_t}.
\end{equation}
Par l'inégalité de Hadamard, $H^{dep}_t\le 0$, avec égalité si et seulement si
$R_t=I$ : \emph{la corrélation détruit toujours de l'entropie}, et $-H^{dep}_t$
est précisément la multi-information gaussienne (corrélation totale) du vecteur
de rendements. Une lecture utile de la même quantité est spectrale : avec
$\lambda_k$ les valeurs propres de $R_t$ et $\tilde\lambda_k=\lambda_k/N$, le
nombre de participation
$N^{\text{eff}}_t=\exp\!\left(-\sum_k\tilde\lambda_k\log\tilde\lambda_k\right)$
est le nombre effectif de modes de risque indépendants ; il vaut $N$ pour $R=I$
et tend vers $1$ lorsqu'un unique mode collectif absorbe la variance.

L'équation~\eqref{eq:ceiling} est le plafond, pas l'entropie. C'est ce
qu'implique la seule matrice de covariance, et la proposition~\ref{prop:maxent}
dit que la vérité se situe en dessous.

\subsection{Forme : le déficit d'entropie et sa composition}

En soustrayant \eqref{eq:deficit} de \eqref{eq:ceiling} on obtient la
décomposition de l'article, trois canaux et une constante :
\begin{equation}
\label{eq:three}
h(p_t) \;=\; \tfrac{N}{2}\log(2\pi e) \;+\; \underbrace{H^{vol}_t}_{\text{échelle}} \;+\; \underbrace{H^{dep}_t}_{\text{dépendance}} \;-\; \underbrace{\DD_t}_{\text{forme}}.
\end{equation}

Le canal de forme, ce sont l'asymétrie et les queues épaisses, et le
développement d'Edgeworth le dit explicitement : pour une variable standardisée
d'asymétrie $S$ et de kurtosis excédentaire $K$ au voisinage de la gaussianité,
\begin{equation}
\label{eq:edgeworth}
J \;\approx\; \frac{S^2}{12} + \frac{K^2}{48}.
\end{equation}
L'équation~\eqref{eq:edgeworth} répond à la question \og où un saut d'entropie
serait-il visible ? \fg{} : dans les moments d'ordre trois et quatre, de manière
additive, avec un canal d'asymétrie et un canal de queue. Elle est aussi la
source du problème de signe que la section~\ref{sec:dynamics} doit résoudre, car
$S$ y entre au carré --- \emph{la néguentropie ne distingue pas un krach d'une
envolée}. Comme estimateur, \eqref{eq:edgeworth} est inutilisable sur des
rendements quotidiens ; la section~\ref{sec:program} documente à quel point et
ce que nous utilisons à la place.

\subsection{L'indice structurel}

Le canal d'échelle est une nuisance pour notre propos : il domine numériquement
le niveau de \eqref{eq:three}, il bouge avec la volatilité, et un changement
d'unités le modifie. Les deux autres canaux sont sans échelle. En les combinant,
et en utilisant la règle de chaînage exacte
$\KL{p}{\prod_i q_i} = \sum_i \KL{p_i}{q_i} + I(x)$ où $I$ est la
multi-information :

\begin{definition}[Indice d'information structurelle]
\label{def:J}
\begin{equation}
\JJ_t \;\equiv\; \DD_t \;-\; H^{dep}_t \;=\; \KL{p_t}{\textstyle\prod_i q_{i,t}} \;\ge\; 0 ,
\end{equation}
la divergence de la loi jointe conditionnelle des rendements à un produit de
marginales gaussiennes de mêmes échelles.
\end{definition}

$\JJ$ est nul si et seulement si les rendements sont indépendants et gaussiens ;
il s'élève quand les actifs se couplent, quand les marginales s'épaississent ou
quand elles s'asymétrisent. Il est mesuré en nats, il est invariant au niveau de
volatilité, et il est additif sur ses deux canaux, de sorte qu'on peut toujours
demander quelle part d'un mouvement vient de la dépendance et quelle part de la
forme. C'est la variable d'état de la section~\ref{sec:dynamics}.

\section{Dynamique de l'information : un théorème, deux postulats}
\label{sec:dynamics}

\subsection{Les sauts vers le bas sont un théorème}

Deux énoncés distincts sous-tendent la moitié descendante de la dynamique
postulée. Les confondre est une erreur facile, et les versions antérieures de ce
travail la commettaient ; nous les séparons donc explicitement.

Rappelons d'abord la définition de l'\emph{information mutuelle} entre le
vecteur de rendements $X$ et un signal $Y$ :
\begin{equation}
\label{eq:mi}
I(X;Y) \;=\; \KL{p_{XY}}{p_X\otimes p_Y}
\;=\; \iint p(x,y)\,\log\frac{p(x,y)}{p(x)p(y)}\,dx\,dy ,
\end{equation}
c'est-à-dire la divergence de la loi jointe au produit des marginales : de
combien la loi jointe s'écarte de l'indépendance. Elle est positive, symétrique,
nulle si et seulement si $X\perp Y$, et --- contrairement à $h$ --- invariante par
reparamétrisation inversible. Pour un couple gaussien scalaire de corrélation
$\rho$, $I=-\tfrac12\log(1-\rho^2)$.

\begin{proposition}[Chute de l'entropie conditionnelle moyenne]
\label{prop:conditioning}
Soit $X$ le vecteur de rendements et $Y$ un signal. Alors
\begin{equation*}
\mathbb{E}_Y\!\left[h(X\mid Y)\right] \;=\; h(X) - I(X;Y) \;\le\; h(X),
\end{equation*}
avec égalité si et seulement si $X\perp Y$.
\end{proposition}

\begin{proposition}[Monotonie du plafond sous contrainte]
\label{prop:constraint}
Soient $\mathcal{C}\subseteq\mathcal{C}'$ deux ensembles de contraintes et
$\mathcal{Q}(\cdot)$ l'ensemble des lois qui les satisfont. Alors
$\mathcal{Q}(\mathcal{C}')\subseteq\mathcal{Q}(\mathcal{C})$ et donc
\begin{equation*}
\sup_{p\,\in\,\mathcal{Q}(\mathcal{C}')} h(p) \;\le\; \sup_{p\,\in\,\mathcal{Q}(\mathcal{C})} h(p).
\end{equation*}
\end{proposition}

Les deux résultats sont classiques \citep{cover2006} et ne sont pas
équivalents ; la démonstration détaillée et un exemple gaussien chiffré figurent
en annexe~\ref{app:mi}. Leurs portées diffèrent sur deux axes.

La proposition~\ref{prop:conditioning} \emph{quantifie} la baisse --- exactement
$I(X;Y)$ nats --- mais seulement \emph{en moyenne} sur les réalisations du
signal. Pour un $y$ particulier, $h(X\mid Y=y)$ peut excéder $h(X)$ : un signal
suffisamment déroutant laisse le marché plus incertain qu'avant. Il faut donc
s'attendre à une composante de saut principalement, et non exclusivement,
unilatérale --- ce que trouve la section~\ref{sec:program}.

La proposition~\ref{prop:constraint} vaut au contraire \emph{pour chaque
réalisation}, mais porte sur le \emph{plafond} de maximum d'entropie et non sur
l'entropie réelle, et elle ne quantifie rien : elle dit que le plafond ne peut
pas monter, pas de combien il descend.

Ce que les deux disent ensemble, et qui motive tout ce qui suit : la chute est
\emph{discontinue}, car le signal arrive en un point du temps, et sa taille
espérée en nats est le contenu informationnel de la nouvelle --- une amplitude
mesurée sur une échelle qui ne dépend ni de l'actif, ni de la devise, ni des
unités.

\paragraph{Une réserve d'honnêteté.} Aucune de ces deux propositions n'est
utilisée pour calculer quoi que ce soit dans la partie empirique. Nous mesurons
$\JJ$, et non $h(X)-\mathbb{E}_Y[h(X\mid Y)]$ : estimer une entropie
différentielle en dimension~20 sur une fenêtre de 252~jours est hors de portée
(annexe~\ref{app:units}). Le lien entre le théorème et les statistiques
rapportées plus loin est un raisonnement, pas une identité estimée, et le
lecteur doit le tenir pour tel.

\subsection{Deux postulats}

\begin{postulate}[Relaxation]
\label{p:relax}
En l'absence d'information nouvelle, l'indice structurel décroît vers un niveau
de long terme à taux constant :
$d\JJ_t = -\kappa\left(\JJ_t-\bar\JJ\right)dt + \sigma_\JJ\,dW_t$ avec $\kappa>0$.
\end{postulate}

L'économie derrière P\ref{p:relax} est qu'une contrainte s'évente. La nouvelle
est intégrée aux prix, le désaccord se reforme, les positions se dispersent, et
la coupe transversale recouvre la liberté que le choc lui avait retirée. Rien
dans la théorie de l'information ne délivre cela ; c'est l'hypothèse substantielle
de l'article, et $\kappa$ est l'objet à estimer.

\paragraph{Pourquoi un processus d'Ornstein--Uhlenbeck et pas autre chose ?}
Le choix n'est pas un choix de commodité : il est \emph{forcé} dès qu'on impose
trois exigences minimales, et il est le seul à les satisfaire toutes les trois.

\begin{enumerate}
\item \emph{Markov.} L'avenir de $\JJ$ ne dépend du passé qu'à travers son
niveau présent. C'est l'exigence de parcimonie : ne postuler aucune mémoire
au-delà de ce que l'état courant contient déjà.
\item \emph{Stationnaire.} $\JJ$ possède un niveau de long terme et une variance
finis. Sans cela l'indice dérive, il n'est pas une variable d'état, et les
comparaisons de régime de la section~\ref{sec:program} n'ont pas de sens.
\item \emph{Bruit gaussien.} Entre deux arrivées d'information, on ne prétend
rien savoir du bruit au-delà de sa variance --- c'est le principe de maximum
d'entropie appliqué au bruit lui-même, exactement comme il l'a été aux
rendements en section~\ref{sec:bridge}.
\end{enumerate}

Le théorème de \citet{doob1942} énonce que \textbf{le processus d'Ornstein--Uhlenbeck
est l'unique processus à la fois markovien, gaussien et stationnaire} en temps
continu. Les trois exigences ne laissent donc aucune latitude sur la partie
diffusive : elle est déterminée. Le terme de saut, lui, est ajouté précisément
pour \emph{briser} la gaussianité là où le cadre affirme qu'elle ne tient pas,
c'est-à-dire aux arrivées d'information. La structure de \eqref{eq:ou} est ainsi
un squelette imposé plus une composante délibérément non gaussienne, ce qui est
exactement la classe introduite par \citet{bns2001}.

\paragraph{Ce qui, en revanche, reste arbitraire.} Trois choses, qu'il faut
énoncer plutôt que laisser croire qu'elles découlent du théorème. Le taux
$\kappa$ est supposé \emph{constant} : rien n'interdit qu'une crise soit
absorbée plus vite, ou plus lentement, qu'une nouvelle ordinaire, et les données
employées ici ne permettent pas de trancher. Le rappel est \emph{linéaire} en
l'écart au niveau de long terme, alors qu'un rappel accéléré loin de la moyenne
serait tout aussi plausible. Enfin, $\JJ$ est positif par construction alors
qu'un Ornstein--Uhlenbeck muni d'un terme brownien peut devenir négatif : la
spécification de \citet{bns2001} évite cet écueil en supprimant le terme
brownien et en pilotant le processus par un subordinateur à sauts positifs, ce
qui garantit une loi stationnaire portée par $\mathbb{R}^+$. Empiriquement la
question ne se pose pas --- $\JJ$ reste compris entre $2{,}17$ et $21{,}13$~nats
sur le panel principal et entre $0{,}42$ et $1{,}69$ sur le panel de
réplication, donc très loin de la frontière --- mais formellement il faudrait
soit supprimer $\sigma_\JJ$, soit écrire le processus sur $\log\JJ$.

Le postulat n'est enfin pas anodin sur le plan de la mesure : un estimateur doté
de mémoire fabriquera une décroissance même en l'absence de toute relaxation
économique, aussi la section~\ref{sec:program} mesure-t-elle la décroissance
qu'un monde sans mémoire produirait déjà.

\begin{postulate}[Asymétrie de la contrainte]
\label{p:sign}
La contrainte qu'impose l'information est unilatérale dans la direction des
pertes, de sorte qu'un saut d'entropie s'accompagne d'une asymétrie négative du
facteur de marché.
\end{postulate}

\paragraph{Pourquoi ce postulat est logiquement nécessaire.} L'indice $\JJ$ est
une fonction \emph{paire} de l'asymétrie : changer le signe de l'asymétrie ne
change pas sa valeur. La raison est immédiate. Le canal de forme est
$\Dun(z)=\KL{p_z}{\varphi}$, et la loi normale $\varphi$ est symétrique ; donc
si l'on remplace $z$ par $-z$, la loi de référence est inchangée et la
divergence l'est aussi :
\begin{equation}
\label{eq:even}
\Dun(-z)=\Dun(z),\qquad\text{et par conséquent}\qquad \JJ(-x)=\JJ(x).
\end{equation}
Le développement d'Edgeworth \eqref{eq:edgeworth} le rend visible à l'\oe il :
l'asymétrie $S$ y entre au carré.

La conséquence est structurelle et non anecdotique. \textbf{Une distribution à
longue queue à gauche --- un krach --- et son image en miroir à longue queue à
droite --- une envolée --- ont exactement le même $\JJ$.} L'indice mesure la
\emph{taille} de l'information, jamais son sens. Aucun raffinement de
l'estimateur n'y changera quoi que ce soit : c'est une propriété de l'objet
mesuré. Si l'article veut dire quoi que ce soit sur la direction, il doit donc
ajouter un ingrédient extérieur. P\ref{p:sign} est cet ingrédient, et c'est le
plus petit possible : une seule affirmation, sur le signe, ajoutée à un cadre
qui par construction n'en produit aucune.

\paragraph{Sur quoi il s'appuie.} Non pas sur un mécanisme que nous
affirmerions, mais sur une régularité empirique établie : l'\emph{effet de
levier} au sens de \citet{black1976}, c'est-à-dire le fait que les rendements
négatifs sont suivis de hausses de volatilité plus fortes que les rendements
positifs de même amplitude. C'est l'un des faits stylisés les plus robustes de
la finance empirique \citep{christie1982,engle1993,nelson1991}. Traduit dans le
langage de cet article : à quantité d'information égale, la réponse du marché
est plus ample vers le bas que vers le haut, donc un saut de $\JJ$ arrive
accompagné d'une asymétrie de marché négative.

Nous nous abstenons délibérément de choisir entre les explications causales
concurrentes de cet effet --- levier financier à proprement parler, rétroaction
de la volatilité, ou retrait asymétrique de la liquidité du côté vendeur. Le
postulat ne dépend d'aucune d'entre elles : il n'affirme que l'asymétrie, pas sa
cause. Il est testable séparément (H5) et réfutable séparément, et la
section~\ref{sec:program} conclut d'ailleurs qu'il n'est pas établi.

\subsection{La dynamique qui en résulte}

La proposition~\ref{prop:conditioning} jointe à P\ref{p:relax} donne un processus
d'Ornstein--Uhlenbeck non gaussien à sauts unilatéraux, de la classe introduite
par \citet{bns2001} pour la volatilité stochastique :
\begin{equation}
\label{eq:ou}
d\JJ_t \;=\; -\kappa\left(\JJ_t-\bar\JJ\right)dt \;+\; \sigma_\JJ\,dW_t \;+\; \int_{z>0} z\,\mathcal{P}(dt,dz),
\end{equation}
Le terme intégral demande d'être lu avec soin, car sa notation est compacte.

\paragraph{Ce que $\mathcal{P}$ n'est pas.} Ce n'est ni une loi gaussienne, ni
une fonction. La lettre $\mathcal N$ est réservée dans cet article à la loi
normale, et la mesure de saut porte donc un symbole distinct.

\paragraph{Ce que $\mathcal{P}$ est.} Une \emph{mesure aléatoire de comptage},
dite mesure aléatoire de Poisson. On peut se la représenter comme un semis
aléatoire de points dans le plan $(\text{date},\ \text{taille})$ : chaque point
est un événement informationnel, son abscisse est la date à laquelle il survient
et son ordonnée la taille du saut qu'il provoque. Formellement,
$\mathcal{P}(dt,dz)$ vaut $1$ s'il existe un saut dans l'intervalle de temps
$[t,t+dt)$ dont la taille tombe dans $[z,z+dz)$, et $0$ sinon. L'intégrale
$\int_{z>0} z\,\mathcal{P}(dt,dz)$ se lit alors simplement : \og s'il y a un
saut à la date $t$, ajouter sa taille $z$ à $\JJ$ ; sinon n'ajouter rien \fg{}.
C'est une écriture condensée de $\sum_{j} z_j$, la somme des sauts survenus.

\paragraph{Ce qui la caractérise.} Deux ingrédients, réunis dans sa
\emph{mesure d'intensité} $\nu(dz)\,dt$, qui est le nombre espéré de sauts de
taille comprise entre $z$ et $z+dz$ pendant $dt$ : d'une part une
\emph{fréquence d'arrivée} $\lambda=\nu(\mathbb{R}^+)$, en nombre de sauts par
unité de temps, d'autre part une \emph{loi des tailles} $\nu(dz)/\lambda$. Sur
le panel principal, la section~\ref{sec:program} estime $\lambda\approx5$ sauts
par an, une taille moyenne de $0{,}62$~nat et un maximum de $3{,}21$~nats.

\paragraph{Ce qu'est $z$.} La taille d'un saut individuel, exprimée en nats.
Conceptuellement, c'est l'information mutuelle de la nouvelle qui l'a provoqué,
au sens de la proposition~\ref{prop:conditioning}. Opérationnellement, c'est la
variation de $\JJ$ détectée par le filtre à quatre sigmas de la
section~\ref{sec:program}. Les deux ne coïncident qu'approximativement, et le
papier ne prétend pas les identifier.

\paragraph{Pourquoi l'intégrale ne porte que sur $z>0$.} Parce que
$\JJ$ mesure un \emph{déficit} d'entropie et que, par la
proposition~\ref{prop:conditioning}, l'information ne peut qu'accroître ce
déficit en espérance. Dans ce cadre, rien ne crée d'\emph{augmentation soudaine}
d'entropie : l'entropie n'est restituée que par la relaxation lente du premier
terme, jamais par un saut. D'où des sauts unilatéraux.

Deux réserves, toutefois. La proposition~\ref{prop:conditioning} est un énoncé
\emph{en moyenne} : une réalisation particulièrement déroutante peut accroître
l'entropie, donc une spécification fidèle autoriserait une petite probabilité de
sauts négatifs plutôt que de les interdire. Et la mesure de Poisson suppose des
arrivées indépendantes, alors que l'étude d'événement de la
section~\ref{sec:program} montre un regroupement net des sauts ; un processus
auto-excitant de type Hawkes serait plus fidèle.

Notons enfin que le coefficient de diffusion est noté $\sigma_\JJ$ et non
$\eta$, cette dernière lettre étant réservée au rayon des scénarios de stress
(section~\ref{sec:stress}). Lu sur l'entropie elle-même plutôt que sur
$\JJ$, le signe s'inverse : l'entropie chute par sauts et remonte par diffusion.
Les prédictions observables sont nettes, et ce sont elles que teste la
section~\ref{sec:program} : la loi de $\Delta\JJ$ est asymétrique à droite ;
les grands mouvements sont à la hausse et les petits à la baisse ; et après un
saut positif l'indice décroît avec une demi-vie $\ln 2/\kappa$.

\subsection{Pourquoi queues épaisses, corrélation et sauts arrivent ensemble}

Le cadre produit une prédiction qu'aucun modèle de volatilité ni de corrélation
ne produit seul. Elle se déduit en quatre pas, sans vocabulaire technique.

\begin{enumerate}
\item Une nouvelle importante concerne \emph{tous} les actifs à la fois : c'est
la même information qui contraint chacun d'eux.
\item Si la nouvelle est mauvaise, tous les actifs baissent donc \emph{le même
jour}.
\item Or ce jour-là apparaît aussi, pris isolément, dans l'historique de chaque
actif : c'est pour lui une journée de perte inhabituellement grande.
\item Accumuler de telles journées déforme la distribution des rendements de
chaque actif : le côté des pertes s'allonge et s'alourdit davantage que le côté
des gains.
\end{enumerate}

Le pas 2 est ce que mesure le canal de dépendance : les actifs bougent ensemble,
la corrélation mesurée monte, $-H^{dep}$ monte. Le pas 4 est ce que mesure le
canal de forme : des pertes extrêmes plus nombreuses et plus grandes que les
gains extrêmes, c'est-à-dire un kurtosis plus élevé et une asymétrie négative,
donc $\DD$ monte. Les deux canaux montent ensemble parce qu'ils décrivent le
même jour vu de deux endroits différents --- la coupe transversale pour l'un, la
distribution de chaque actif pour l'autre.

En marché calme, les mouvements marquants sont propres à chaque actif : ils
n'alignent pas la coupe transversale et n'ont donc aucune raison de faire monter
les deux canaux en même temps.

La forme testable est un \emph{couplage dépendant du régime} : la corrélation
entre $-\Delta H^{dep}$ et $\Delta\DD$ devrait être proche de zéro en marché
calme et fortement positive sous stress. C'est H6, et c'est là que se trouve la
réponse de l'article à la question \og le lien entre queues épaisses, corrélation
et sauts d'entropie exige-t-il un postulat supplémentaire ? \fg{} : non. Étant
donné la proposition~\ref{prop:conditioning}, le couplage découle du caractère
commun de la nouvelle. Seul le \emph{signe} exige P\ref{p:sign}.


\section{Le programme empirique H1--H8 et ce que les témoins en laissent}
\label{sec:program}

Cette section résume, sans les reproduire, les estimateurs et les tests de la
version arXiv~v1. Elle sert à une chose : expliquer pourquoi, de toute la
théorie des sections~\ref{sec:bridge} à~\ref{sec:dynamics}, seul le canal de
dépendance passe dans le papier principal.

\subsection{Les huit hypothèses}

\begin{description}
\item[H1 (signature de régime).] Sous stress, $H^{dep}$ baisse, $\DD$ monte, et
donc $\JJ$ monte.
\item[H2 (compensation partielle).] $\Delta H^{dep}_t=-\beta\,\Delta H^{vol}_t+u_t$
avec $\beta>0$, et $\beta_{\text{stress}}>\beta_{\text{calme}}$.
\item[H3 (asymétrie des sauts).] Les variations de $\JJ$ sont unilatérales,
au-delà de ce que produit le même estimateur sur des données sans structure
temporelle.
\item[H4 (relaxation).] $\JJ$ revient à la moyenne avec une demi-vie plus longue
que la mémoire de l'estimateur.
\item[H5 (signature d'asymétrie).] Les sauts positifs de $\JJ$ coïncident avec
une asymétrie négative du facteur de marché (postulat~P\ref{p:sign}).
\item[H6 (couplage des canaux).] Dépendance et forme sont faiblement couplées en
marché calme, fortement sous stress.
\item[H7 (pouvoir prédictif).] L'entropie d'aujourd'hui améliore, hors
échantillon, la prévision du risque à 21~jours par rapport à la volatilité.
\item[H8 (composition du budget de crise).] Regroupés en épisodes, les sauts ont
plus souvent le canal de dépendance pour contributeur dominant qu'un
rééchantillonnage sans mémoire.
\end{description}

\paragraph{Estimation, en une phrase.} Fenêtres glissantes de 252~jours,
covariance rétrécie à la \citet{ledoit2004}, déficit marginal par
l'approximation de \citet{hyvarinen1998} plutôt que par le développement
d'Edgeworth~\eqref{eq:edgeworth} (inutilisable sur des rendements quotidiens),
régime de stress défini par un seuil de volatilité calculé sur le passé
seulement, sauts détectés par un filtre à quatre sigmas sur $\Delta\JJ$.

\subsection{Pourquoi un témoin calibré est indispensable}

Le point théorique qui gouverne toute la lecture des résultats est le suivant.
$\JJ$ est une fonctionnelle \emph{convexe et positive} des moments d'ordre deux à
quatre estimés sur une fenêtre. Quand une observation à queue épaisse entre dans
la fenêtre, les moments estimés sautent, et la fonctionnelle convexe saute
\emph{vers le haut} quel que soit le signe de l'observation. Quand
l'observation vieillit (pondération exponentielle) ou sort de la fenêtre, la
fonctionnelle redescend graduellement. L'appareil de mesure fabrique donc, à
lui seul, la dent de scie que prédisent la proposition~\ref{prop:conditioning}
et le postulat~P\ref{p:relax}. Une simulation gaussienne i.i.d., sans aucun
événement, produit déjà une demi-vie apparente de 35~jours et 37~sauts positifs
contre aucun négatif au filtre à quatre sigmas.

D'où deux témoins, construits en rééchantillonnant la matrice de rendements
observée et en faisant tourner la chaîne identique deux cents fois : le témoin
\emph{i.i.d.} (lignes tirées indépendamment, queues marginales et dépendance
transversale conservées, toute chronologie détruite) et le témoin \emph{par
blocs} de 21~jours (qui conserve en outre le regroupement de volatilité). Une
statistique qui franchit les deux est une preuve ; une statistique qui en
franchit un seul n'est pas identifiée.

\subsection{Le verdict}

\begin{table}[htbp]
\centering
\small
\begin{tabular}{@{}llrrrl@{}}
\toprule
& Statistique & Données & i.i.d. & blocs-21 & Verdict \\
\midrule
H1 & écart de régime $H^{dep}$ & $-0{,}99$ & $-0{,}21$ & $-0{,}51$ & \textbf{franchit les deux} \\
H1 & écart de régime $\JJ$ & $1{,}93$ & $2{,}14$ & $1{,}97$ & échoue aux deux \\
H1 & écart de régime $\DD$ & $0{,}27$ & $1{,}25$ & $0{,}65$ & à contre-sens \\
H2 & $\beta$ (stress) & $0{,}20$ & $0{,}05$ & $0{,}11$ & échoue aux deux \\
H3 & sauts positifs & $158$ & $148{,}5$ & $133{,}2$ & blocs seulement \\
H3 & asymétrie de $\Delta\JJ$ & $12{,}26$ & $14{,}80$ & $12{,}97$ & à contre-sens \\
H4 & demi-vie (jours) & $120{,}8$ & $92{,}1$ & $162{,}8$ & entre les deux \\
H4 & étendue de $\JJ$ & $18{,}95$ & $10{,}05$ & $13{,}77$ & \textbf{franchit les deux} \\
H5 & $\mathrm{corr}(\Delta\JJ,\Delta S_{mkt})$ & $-0{,}24$ & $-0{,}04$ & $-0{,}07$ & échoue aux deux \\
H6 & écart de couplage & $0{,}59$ & $0{,}77$ & $0{,}67$ & à contre-sens \\
H7 & gain hors échantillon & \multicolumn{3}{c}{aucun sur deux panels} & rejetée \\
H8 & part d'épisodes à dominante dépendance & $46\,\%$ & $19$--$33\,\%$ & n.i. & i.i.d.\ seulement \\
\bottomrule
\end{tabular}
\caption{Verdicts face aux témoins (moyennes de deux cents réplications),
panel principal de vingt titres du S\&P~500, 1990--2022. Source :
\texttt{paper/tables/tab\_placebo\_sp500.tex} et sections H7--H8 de la v1.
n.i.\ : non identifié, à toute longueur de bloc testée entre 5 et 126~jours.}
\label{tab:program}
\end{table}

Deux statistiques franchissent les deux témoins sur le panel principal, et une
seule se reproduit sur les 48~portefeuilles sectoriels de Kenneth French :
l'écart de régime du canal de dépendance ($-3{,}65$~nats contre $-0{,}56$ et
$-2{,}46$). La phase~A l'a ensuite durci : il tient en version dévolatilisée
(Forbes--Rigobon) et en version en rangs, cette dernière étant la plus forte
face au témoin par blocs ($p=0{,}03$ sur le S\&P~500, $0{,}01$ sur les secteurs),
alors que le ratio d'absorption de \citet{kritzman2011} échoue au témoin par
blocs sur les deux panels.

\subsection{Ce que le verdict dit de la théorie}

Il faut séparer ce que le verdict touche et ce qu'il ne touche pas.

\paragraph{Ce qui n'est pas touché.} Les énoncés des
sections~\ref{sec:bridge} et~\ref{sec:channels} sont des identités et des
théorèmes. La proposition~\ref{prop:conditioning} aussi. Aucun résultat empirique
ne peut les réfuter ; il peut seulement montrer que les estimateurs disponibles
ne permettent pas de les voir à l'\oe uvre.

\paragraph{Ce qui est touché.} Les deux postulats. P\ref{p:relax} n'est soutenu
que par l'amplitude (l'étendue de $\JJ$), pas par un taux de relaxation : la
demi-vie se situe entre les deux témoins et son niveau dépend du témoin
choisi. P\ref{p:sign} n'est pas établi (H5). Le couplage des canaux (H6), que la
section~\ref{sec:dynamics} dérivait du caractère commun de la nouvelle, est
reproduit et même dépassé par des données rééchantillonnées.

\paragraph{Pourquoi le canal de dépendance survit.} Une explication
vraisemblable, cohérente avec la phase~A, tient à la forme de l'estimateur.
$\tfrac12\log\det R=\tfrac12\sum_k\log\lambda_k$ pèse tout le spectre de la
matrice de corrélation, y compris les petites valeurs propres, qui s'effondrent
quand la coupe transversale se couple sur quelques modes collectifs. Un seul
gros rendement isolé déplace peu ces petites valeurs propres. À l'inverse, le
canal de forme $\DD$ réagit fortement à une seule observation extrême, ce qui
est exactement la signature que reproduit le témoin. Les mesures de
concentration (ratio d'absorption, part de la première valeur propre) ne
regardent que le haut du spectre, où le témoin par blocs les rattrape.

\paragraph{Conséquence pour le papier principal.} L'indicateur retient le seul
canal identifié, en abandonnant $\DD$ et donc $\JJ$. La théorie des
sections~\ref{sec:channels} et~\ref{sec:dynamics} reste la justification de
\emph{pourquoi} ce canal mesure une information partagée (la
multi-information gaussienne, équation~\eqref{eq:ceiling} et annexe~\ref{app:gauss}),
mais la dynamique postulée n'est pas utilisée.

\section{Géométrie des scénarios en boule de Kullback--Leibler}
\label{sec:stress}

Le papier principal retient l'usage le plus simple de cette géométrie : tarifer
en nats un scénario publié, ce qui pour un pur déplacement de moyenne revient à
la moitié du carré de la distance de Mahalanobis. Cette section donne la
construction complète dont cet usage est un cas particulier, avec les
démonstrations que la v1 se contentait d'énoncer, et un résultat nouveau sur le
relèvement d'un scénario de portefeuille à la dimension pleine.

\subsection{Divergence entre deux gaussiennes}

\begin{proposition}[Forme fermée]
\label{prop:klgauss}
Pour $p_1=\mathcal N(\mu_1,\Sigma_1)$ et $p_0=\mathcal N(\mu_0,\Sigma_0)$ sur
$\mathbb R^N$, avec $\delta=\mu_1-\mu_0$,
\begin{equation}
\label{eq:klgauss}
\KL{p_1}{p_0}=\tfrac12\Big[\underbrace{\delta^\top\Sigma_0^{-1}\delta}_{\text{déplacement}}
+\underbrace{\operatorname{tr}(\Sigma_0^{-1}\Sigma_1)-N-\log\det(\Sigma_0^{-1}\Sigma_1)}_{\text{déformation}}\Big].
\end{equation}
\end{proposition}

\begin{proof}
$\KL{p_1}{p_0}=\mathbb E_{p_1}[\log p_1-\log p_0]$. Les deux logarithmes sont
quadratiques ; l'astuce de la trace de l'annexe~\ref{app:gauss} donne
$\mathbb E_{p_1}[(X-\mu_1)^\top\Sigma_1^{-1}(X-\mu_1)]=N$ et
$\mathbb E_{p_1}[(X-\mu_0)^\top\Sigma_0^{-1}(X-\mu_0)]
=\operatorname{tr}(\Sigma_0^{-1}\Sigma_1)+\delta^\top\Sigma_0^{-1}\delta$. On
regroupe avec les termes en $\log\det$.
\end{proof}

Le terme de déplacement est la moitié du carré de la distance de Mahalanobis :
c'est l'unité dans laquelle le papier principal audite les scénarios
publiés. Le terme de déformation est positif et nul si et seulement si
$\Sigma_1=\Sigma_0$ ; en écrivant $\nu_k$ les valeurs propres de
$\Sigma_0^{-1}\Sigma_1$, il vaut $\tfrac12\sum_k(\nu_k-1-\log\nu_k)$, une somme de
termes positifs, un par direction propre.

\subsection{La boule projetée et le taux de change}

\begin{definition}[Scénario en boule d'entropie]
Étant donné la loi conditionnelle d'aujourd'hui $p_0$ et une sévérité $\eta>0$,
le scénario de sévérité $\eta$ pour une fonctionnelle de risque $\rho$ est
\begin{equation}
\label{eq:ball}
p^\star_\eta=\arg\max_{p\,:\,\KL{p}{p_0}\le\eta}\rho(p).
\end{equation}
\end{definition}

Pour un portefeuille linéaire $w^\top r$ de moyenne $m_0$ et d'écart type $s_0$
sous $p_0$, posons $u=(m-m_0)/s_0$ et $v=s/s_0$. Par la
proposition~\ref{prop:klgauss} en dimension un,
\begin{equation}
\label{eq:g}
g(u,v)=\tfrac12\left(u^2+v^2-1-2\log v\right).
\end{equation}

\begin{proposition}[Solution de la boule projetée]
\label{prop:ballsol}
Pour $\rho=-m+c\,s$ avec $c\ge0$ (VaR et expected shortfall gaussiennes sont de
cette forme), le maximum de $\rho$ sous $g(u,v)\le\eta$ est atteint sur la
frontière $g=\eta$ en un point vérifiant
\begin{equation*}
u=-t,\qquad v-v^{-1}=c\,t,\qquad t>0 .
\end{equation*}
Pour $c=0$ : $v=1$ et $u=-\sqrt{2\eta}$. \emph{Un scénario purement
directionnel de $\eta$~nats est un mouvement de $\sqrt{2\eta}$~sigmas} ; un
nat achète $\sqrt2\approx1{,}41$~sigma, et $k$~sigmas coûtent $k^2/2$~nats.
\end{proposition}

\begin{proof}
$g$ est convexe sur $\mathbb R\times(0,\infty)$ (somme de $u^2/2$ et de
$(v^2-1-2\log v)/2$, de dérivée seconde $1+v^{-2}>0$) et l'objectif
$s_0(-u+cv)$ est linéaire : les conditions de Karush--Kuhn--Tucker sont
suffisantes. L'objectif n'a pas de maximum intérieur, donc la contrainte est
saturée. Le lagrangien donne $-1=\lambda u$ et $c=\lambda(v-v^{-1})$ avec
$\lambda>0$ ; en posant $t=1/\lambda$ on obtient le système annoncé, et $t$ est
fixé par $g(-t,v(t))=\eta$. Pour $c=0$, $v-v^{-1}=0$ impose $v=1$ et
$g=u^2/2=\eta$.
\end{proof}

\subsection{Du portefeuille à la dimension pleine}

La boule~\eqref{eq:ball} se résout sur la loi du portefeuille, mais un test de
résistance doit dire comment chaque actif bouge. Deux faits règlent la
question.

\begin{proposition}[Traitement des données]
\label{prop:dpi}
Pour toute loi $p$ sur $\mathbb R^N$, et $p^w$, $p_0^w$ les lois de $w^\top r$
sous $p$ et $p_0$,
$\KL{p^w}{p_0^w}\le\KL{p}{p_0}$.
\end{proposition}

C'est l'inégalité de traitement des données \citep{cover2006} appliquée à
l'application $r\mapsto w^\top r$. Le prix \emph{ambiant} d'un scénario (en
dimension pleine) est donc toujours au moins son prix \emph{vu du portefeuille}
(en projection). Le tableau des bundles classiques de la v1 le vérifie ligne à
ligne ($9{,}97$ contre $7{,}24$~nats pour le bundle standard).

\begin{proposition}[Relèvement minimal]
\label{prop:lift}
Soit $p_0=\mathcal N(\mu_0,\Sigma_0)$, $m_0=w^\top\mu_0$, $s_0^2=w^\top\Sigma_0w$,
et une cible $(m,s)$ avec $s>0$. Parmi les gaussiennes
$\mathcal N(\mu_1,\Sigma_1)$ telles que $w^\top\mu_1=m$ et
$w^\top\Sigma_1w=s^2$, la divergence $\KL{p_1}{p_0}$ est minimale pour
\begin{equation}
\label{eq:lift}
\mu_1=\mu_0+u\,s_0\,\beta,\qquad
\Sigma_1=\Sigma_0+(v^2-1)\,s_0^2\,\beta\beta^\top,
\qquad\beta=\frac{\Sigma_0w}{s_0^2},
\end{equation}
et ce minimum vaut exactement $g(u,v)$, le prix vu du portefeuille.
\end{proposition}

\begin{proof}
Par la proposition~\ref{prop:dpi}, tout candidat coûte au moins le prix de sa
projection, qui est $g(u,v)$ puisque la projection d'une gaussienne est
gaussienne de moyenne $m$ et d'écart type $s$. Il suffit donc de vérifier que
\eqref{eq:lift} atteint cette borne. On a $w^\top\beta=1$, donc
$w^\top\mu_1=m_0+us_0=m$ et $w^\top\Sigma_1w=s_0^2+(v^2-1)s_0^2=s^2$. Dans
\eqref{eq:klgauss} : $\delta^\top\Sigma_0^{-1}\delta=u^2s_0^2\,\beta^\top\Sigma_0^{-1}\beta=u^2$ ;
$\operatorname{tr}(\Sigma_0^{-1}\Sigma_1)=N+(v^2-1)s_0^2\,\beta^\top\Sigma_0^{-1}\beta=N+v^2-1$ ;
et par le lemme du déterminant matriciel,
$\det(\Sigma_0^{-1}\Sigma_1)=1+(v^2-1)s_0^2\,\beta^\top\Sigma_0^{-1}\beta=v^2$.
D'où $\KL{p_1}{p_0}=\tfrac12(u^2+v^2-1-2\log v)=g(u,v)$. La matrice $\Sigma_1$ est
définie positive pour tout $v>0$.
\end{proof}

\paragraph{Lecture actif par actif.} Le vecteur $\beta$ est celui des bêtas des
actifs sur le portefeuille, $\beta_i=\operatorname{Cov}(r_i,w^\top r)/s_0^2$.
Écrivons $r=\beta f+e$ avec $f=w^\top r$ et $e$ non corrélé à $f$ sous $p_0$.
Le scénario~\eqref{eq:lift} multiplie l'écart type du facteur portefeuille $f$
par $v$, déplace sa moyenne de $u$~sigmas, et laisse la loi de $e$ inchangée.
Chaque actif voit sa variance augmenter de $(v^2-1)s_0^2\beta_i^2$ et sa
moyenne bouger de $us_0\beta_i$ : les actifs à fort bêta prennent le choc, et
les corrélations entre eux montent. C'est un scénario à un facteur, et il est
le moins informatif de tous ceux qui atteignent la sévérité visée.

\begin{remark}[Correction de la v1]
La version arXiv~v1 identifiait ce minimum par ses conditions du premier ordre,
$\Sigma_1^{-1}=\Sigma_0^{-1}-\lambda ww^\top$, et le jugeait \og non
interprétable actif par actif \fg{}. Par la formule de Sherman--Morrison, cette
perturbation de rang un de la précision \emph{est}~\eqref{eq:lift}, une
perturbation de rang un de la covariance le long de $\Sigma_0w$, et elle
s'interprète comme ci-dessus. La v1 ne l'utilisant nulle part, aucun chiffre
n'en dépend.
\end{remark}

\begin{corollary}[Coût de la dilatation uniforme]
\label{cor:uniform}
La dilatation uniforme $\Sigma_1=v^2\Sigma_0$, choix par défaut de la v1, a une
composante de déformation égale à $N\cdot\tfrac12(v^2-1-2\log v)$, soit $N$~fois
celle du relèvement minimal.
\end{corollary}

\begin{proof}
Toutes les valeurs propres de $\Sigma_0^{-1}\Sigma_1$ valent $v^2$ ; on applique
la forme spectrale du terme de déformation.
\end{proof}

Pour le scénario décennal du panel principal ($v=1{,}63$, $N=20$), la
composante de variance vaut $0{,}34$~nat au minimum et $20\times0{,}34\approx6{,}8$
sous la dilatation uniforme, ce qui est le chiffre du tableau de composition de
la v1. Le corollaire explique le constat empirique de la v1 selon lequel les
compositions calibrées sur H1 ($2{,}88$~nats) et raffinées par H8 ($0{,}71$)
coûtent \emph{moins} que le défaut : le défaut déforme également toutes les
directions propres, alors qu'une composition qui concentre le choc sur le mode
commun se rapproche du scénario à un facteur~\eqref{eq:lift}, dont le coût total
est $0{,}41$~nat avec le déplacement de moyenne. Que la dilatation uniforme
coûte $N$~fois le minimum en variance est un résultat général ; que les
compositions calibrées se placent entre les deux reste un constat sur ces
panels, mais il devient attendu plutôt que surprenant.

\subsection{Composer une sévérité donnée}

\begin{proposition}[Une sévérité, plusieurs compositions]
\label{prop:composition}
Si $w^\top\Sigma_aw=w^\top\Sigma_bw=s^2$, alors pour tout $\alpha\in[0,1]$,
$\Sigma_\alpha=(1-\alpha)\Sigma_a+\alpha\Sigma_b$ vérifie
$w^\top\Sigma_\alpha w=s^2$.
\end{proposition}

\begin{proof}
$w^\top\Sigma w$ est linéaire en $\Sigma$.
\end{proof}

La sévérité du portefeuille ne contraint donc pas la composition. Le prix
ambiant, lui, est convexe en $\Sigma_1$ (le terme $-\log\det$ est convexe et le
reste est linéaire), donc le long d'un mélange il est majoré par le mélange des
prix, ce qui garantit qu'une composition intermédiaire ne coûte jamais plus que
la plus chère de ses deux extrémités.

\paragraph{Un plafond physique pour la corrélation.} À volatilités marginales
fixées, $w^\top\Sigma w$ est maximal sur les matrices de corrélation valides pour
$\rho_{ij}=1$, où il vaut $(\sum_iw_i\sigma_i)^2$. Sur le panel principal ce
plafond ne vaut qu'environ $1{,}5$~fois la volatilité d'aujourd'hui : au-delà de
l'échelon quinquennal, aucune hausse de corrélation seule ne suffit, et une part
de hausse des volatilités est inévitable.

\subsection{Calibrer le rayon, et la limite de cette calibration}

Le rayon se lit sur le flux d'information de la projection,
$I^{(h)}_t=\KL{p^w_t}{p^w_{t-h}}$, dont les quantiles (sur blocs non
chevauchants de 21~jours) convertissent une période de retour en rayon. Lire le
rayon sur le flux à $N$ dimensions le surestime d'un ordre de grandeur, par la
proposition~\ref{prop:dpi}. Calibrée ainsi, l'échelle de la v1 donne un
multiplicateur de volatilité décennal de $1{,}63$, quand le quantile décennal de
la hausse de volatilité réalisée sur un mois vaut $3{,}2$. L'échelle est une
règle cohérente et reproductible ; ses périodes de retour ne sont pas validées.
C'est l'une des raisons pour lesquelles le papier principal garde l'audit
(tarifer un scénario donné) et abandonne la génération (construire un scénario
d'une période de retour donnée).

\subsection{Le test de résistance inversé}

Demander le scénario le moins coûteux qui produit une perte $L$ donnée inverse
la boule. Dans le cadre gaussien projeté, la
proposition~\ref{prop:ballsol} se lit à l'envers. Dans le cadre non
paramétrique, la solution est l'inclinaison exponentielle de l'échantillon
\citep{csiszar1975}, dont la taille d'échantillon effective s'effondre quand le
rayon croît : à $9$~nats, elle vaut $1{,}0$ sur le panel principal, c'est-à-dire
que la loi inclinée n'est plus que la pire journée observée. Toute
implémentation devrait rapporter cette taille effective et basculer vers la
boule gaussienne dès qu'elle tombe sous quelques dizaines.

\section{Le pont vers le papier principal : décomposition par blocs}
\label{sec:blocks}

Cette section relie la théorie qui précède à l'indicateur du papier principal.
Elle établit l'identité sur laquelle repose la \og carte de contagion \fg{} et
montre où le canal de dépendance se place dans l'indice structurel.

\subsection{L'identité}

Partitionnons les $N$ séries en $B$ blocs (par exemple les segments du CISS :
monétaire, obligataire, actions, intermédiaires financiers, change), et notons
$R_{bb}$ les blocs diagonaux de la matrice de corrélation $R$ et
$R_D=\operatorname{diag}(R_{11},\dots,R_{BB})$.

\begin{proposition}[Décomposition intra/inter-blocs]
\label{prop:blocks}
\begin{equation}
\label{eq:blocks}
\underbrace{-\tfrac12\log\det R}_{\text{dépendance totale}}
\;=\;\sum_{b=1}^{B}\underbrace{\left(-\tfrac12\log\det R_{bb}\right)}_{\text{intra-bloc }b}
\;+\;\underbrace{\left(-\tfrac12\log\det\big(R_D^{-1/2}RR_D^{-1/2}\big)\right)}_{\text{inter-blocs}} ,
\end{equation}
et chacun des $B+1$ termes est positif ou nul. Le terme inter-blocs est la
multi-information gaussienne entre les blocs,
$\KL{\mathcal N(0,R)}{\bigotimes_b\mathcal N(0,R_{bb})}$ ; il est nul si et
seulement si les blocs sont deux à deux non corrélés.
\end{proposition}

\begin{proof}
$R=R_D^{1/2}\,\big(R_D^{-1/2}RR_D^{-1/2}\big)\,R_D^{1/2}$ et le déterminant est
multiplicatif, donc $\log\det R=\log\det R_D+\log\det(R_D^{-1/2}RR_D^{-1/2})$, avec
$\log\det R_D=\sum_b\log\det R_{bb}$. Chaque $R_{bb}$ est une matrice de
corrélation, donc $-\tfrac12\log\det R_{bb}\ge0$ par l'inégalité de Hadamard
(annexe~\ref{app:gauss}). La matrice $\tilde R=R_D^{-1/2}RR_D^{-1/2}$ a des blocs
diagonaux égaux à l'identité, donc une diagonale unité, et Hadamard s'applique
encore. Enfin la proposition~\ref{prop:klgauss} avec $\mu_1=\mu_0$,
$\Sigma_1=R$, $\Sigma_0=R_D$ donne
$\tfrac12[\operatorname{tr}(R_D^{-1}R)-N-\log\det(R_D^{-1}R)]$, et
$\operatorname{tr}(R_D^{-1}R)=\sum_b\operatorname{tr}(I)=N$.
\end{proof}

\paragraph{Deux blocs : les corrélations canoniques.} Pour $B=2$, en notant
$C=R_{11}^{-1/2}R_{12}R_{22}^{-1/2}$ et $\rho_1,\dots,\rho_k$ ses valeurs
singulières (les corrélations canoniques entre les deux blocs),
$\det\tilde R=\det(I-CC^\top)=\prod_j(1-\rho_j^2)$ et le terme inter-blocs vaut
$-\tfrac12\sum_j\log(1-\rho_j^2)$. C'est l'information mutuelle gaussienne
entre les deux segments, somme des informations portées par chaque paire
canonique. Pour $B>2$, l'identité~\eqref{eq:blocks} s'applique de façon
emboîtée et permet d'attribuer la contagion à des paires ou des groupes de
segments par règle de chaînage.

\subsection{Où se place le canal de dépendance dans l'indice structurel}

La règle de chaînage de l'annexe~\ref{app:index} écrit
$\JJ=\sum_i\Dun(x_i)+I(x)$. En remplaçant la multi-information $I(x)$ par sa
part gaussienne $-H^{dep}$ et en appliquant~\eqref{eq:blocks} :
\begin{equation}
\label{eq:jblocks}
\JJ\;\approx\;\underbrace{\sum_i\Dun(x_i)}_{\text{forme}}
\;+\;\underbrace{\sum_b\left(-\tfrac12\log\det R_{bb}\right)}_{\text{intra}}
\;+\;\underbrace{\left(-\tfrac12\log\det\tilde R\right)}_{\text{inter}} .
\end{equation}
Le papier principal conserve les deux derniers termes et abandonne le premier,
qui est celui que le témoin reproduit (section~\ref{sec:program}). Le CISS
entropique est donc l'indice structurel amputé de son canal de forme, et sa
décomposition est exacte, sans approximation, parce qu'elle ne porte que sur la
partie gaussienne.

\subsection{Trois propriétés utiles à l'indicateur}

\paragraph{Invariance d'échelle.} Les termes de~\eqref{eq:blocks} ne dépendent
que de $R$ : ils sont invariants à toute remise à l'échelle série par série,
ce qui permet d'agréger des marchés aux volatilités très différentes sans
pondération ad hoc. La dévolatilisation de \citet{forbes2002} n'agit donc pas
sur la formule mais sur le biais de l'\emph{estimateur} de $R$ en période
d'hétéroscédasticité.

\paragraph{Version en rangs.} En remplaçant $R$ par la corrélation des scores
normaux $\Phi^{-1}(\hat F_i(x_i))$, on estime la multi-information d'une copule
gaussienne. La quantité devient invariante à toute transformation monotone des
marges, ce qui explique sa robustesse aux queues observée en phase~A.

\paragraph{Normalisation en $N$.} Pour une corrélation constante $\rho$,
$\det R=(1-\rho)^{N-1}\big(1+(N-1)\rho\big)$, de sorte que $-H^{dep}$ croît à peu
près linéairement en $N$. Diviser par $N-1$, ou rapporter la valeur à la
corrélation constante équivalente $\rho^\star$ qui donne le même déterminant,
rend l'indicateur comparable d'un panel à l'autre. Pour $\rho^\star$ petit,
$-H^{dep}\approx\tfrac14N(N-1)\rho^{\star2}$.

\section{Indicateurs de suivi et solvabilité en assurance}
\label{sec:monitoring}

\paragraph{Deux indicateurs de la v1.} La séparation saut/diffusion livrait deux
quantités quotidiennes : la composante de saut positif standardisée de $\JJ$
(quelle quantité d'information est arrivée aujourd'hui) et $\JJ_t$ rapporté à
son maximum glissant (quelle part n'est pas encore résorbée). Ni l'un ni l'autre
n'améliore la prévision de la volatilité (H7), et tous deux reposent sur des
statistiques que le témoin n'identifie pas. Ils ne sont pas repris.

\paragraph{Assurance.} Trois liens restent valables en principe. Un scénario en
boule d'entropie produit une loi stressée complète, objet naturel pour l'ORSA.
La décomposition des bundles classiques concerne la formule standard de
Solvabilité~II, dont le crédit de diversification repose sur une matrice de
corrélation fixe. Enfin les mesures de risque entropiques sont cohérentes
\citep{ahmadijavid2012,artzner1999}. Rien n'a cependant été testé sur des
passifs, et les panels courts et peu fréquents des données d'assurance mettraient
à mal chacun des estimateurs employés.

\section{Questions ouvertes}
\label{sec:open}

\begin{enumerate}
\item \textbf{Arrivées auto-excitantes.} Le regroupement des sauts observé dans
l'étude d'événement de la v1 contredit la mesure de Poisson de~\eqref{eq:ou}.
Un processus de Hawkes est la spécification suivante.
\item \textbf{Positivité.} $\JJ\ge0$ alors qu'un Ornstein--Uhlenbeck à terme
brownien peut devenir négatif. Il faudrait soit supprimer $\sigma_\JJ$ (forme
pure de \citealt{bns2001}), soit modéliser $\log\JJ$.
\item \textbf{Dépendance de queue non linéaire.} $\JJ$ et le canal de dépendance
ne voient que la part gaussienne de la multi-information. Une estimation de la
multi-information de copule complète (par exemple par plus proches voisins sur
les rangs) mesurerait ce qui manque, au prix d'une variance d'estimation bien
plus forte.
\item \textbf{Queue du flux d'information.} Au-delà d'une vingtaine d'années,
les quantiles de $I^{(h)}$ sont extrapolés. Un ajustement de Pareto généralisée
serait la correction disciplinée.
\item \textbf{Relèvement minimal et scénarios publiés.} La
proposition~\ref{prop:lift} donne un étalon : pour tout scénario publié, le
rapport entre son prix ambiant et le prix minimal $g(u,v)$ de sa projection
mesure combien d'information il dépense hors de la direction qui compte pour le
portefeuille. C'est un prolongement direct de l'audit du papier principal.
\end{enumerate}
\appendix

\section{Unités, invariance, et ce qui est calculable}
\label{app:units}

\paragraph{Le nat.} L'entropie $H=-\sum_i p_i\log p_i$ n'est définie qu'une fois
la base du logarithme choisie, et ce choix \emph{est} le choix de l'unité. Avec
$\log_2$ on obtient des bits (shannons), avec $\log_{10}$ des dits, et avec le
logarithme népérien des \emph{nats}. Un nat est l'information apportée par une
observation qui divise l'espace des possibles par $e\approx2{,}718$, là où un bit
le divise par deux. La conversion est
$1~\text{nat}=1/\log 2\approx1{,}4427$~bit. Tout ce papier emploie le logarithme
népérien : entropies, déficits, divergences, rayons de scénario et indice
structurel sont en nats.

Trois lectures opérationnelles, de la plus utile à la plus abstraite. D'abord le
taux de change établi en section~\ref{sec:stress} : un scénario purement
directionnel de $\eta$ nats est un mouvement de $\sqrt{2\eta}$ sigmas, donc
$k$~sigmas coûtent $k^2/2$ nats. Ensuite le codage : $\KL{p}{q}$ nats est
l'excès de longueur de code attendu lorsqu'on encode des données de loi $p$ avec
un code optimal pour $q$. Enfin le test d'hypothèse : une divergence de un nat
par observation signifie que le rapport de vraisemblance en faveur de $p$ croît
en moyenne d'un facteur $e$ à chaque observation.

\paragraph{$h$ dépend des unités, les divergences non.} L'entropie différentielle
d'une variable continue vérifie, pour toute transformation affine inversible,
\begin{equation}
\label{eq:h-scaling}
h(aX+b) \;=\; h(X) + \log|a| ,
\qquad\text{et en dimension } N,\quad h(AX+b)=h(X)+\log|\det A| .
\end{equation}
Exprimer les rendements en pourcentages plutôt qu'en décimales ajoute donc
$N\log 100$ à $h$ : le \emph{niveau} de $h$, et par conséquent celui de
$H^{vol}$ et de $H^{cov}$, n'a aucune valeur absolue interprétable. En revanche
$\KL{p}{q}$, $I(X;Y)$, $\DD$ et $\JJ$ sont invariants par tout changement de
variable inversible, les jacobiens se simplifiant entre numérateur et
dénominateur du logarithme. \emph{Un chiffre en nats n'est comparable d'un jeu
de données à l'autre que si c'est une divergence.} C'est la raison structurelle
pour laquelle l'article construit un indice sans échelle.

\paragraph{Comment calcule-t-on $h(X)$ ?} Dans le cas gaussien, en forme fermée
(annexe~\ref{app:gauss}). Dans le cas général, on ne le peut pas sans connaître
$p$. Des estimateurs non paramétriques existent --- espacements de
\citet{vasicek1976}, plus proches voisins de Kozachenko--Leonenko, méthodes à
noyau --- mais tous se dégradent rapidement au-delà de trois ou quatre
dimensions. Estimer une entropie différentielle en dimension~20 sur une fenêtre
de 252~jours est hors de portée.

L'article ne calcule donc \emph{jamais} $h(p)$. Il calcule $h(q)$ en forme
fermée à partir de $\widehat\Sigma$, et le déficit $\DD$ par une approximation
de la \emph{différence} d'entropies marginales. Une différence d'entropies est
sans dimension et estimable ; une entropie différentielle nue ne l'est pas.
Toute l'architecture --- décomposer plutôt qu'estimer directement --- découle de
cette contrainte.

\section{Entropie gaussienne : démonstration de l'équation~\eqref{eq:ceiling}}
\label{app:gauss}

\begin{proof}[Étape A : entropie d'une gaussienne multivariée]
Pour $q=\mathcal N(\mu,\Sigma)$ sur $\mathbb{R}^N$,
\[
\log q(x) = -\tfrac12\left[N\log 2\pi + \log\det\Sigma + (x-\mu)'\Sigma^{-1}(x-\mu)\right],
\]
donc
\[
h(q) = -\mathbb{E}_q[\log q] = \tfrac12 N\log 2\pi + \tfrac12\log\det\Sigma
+ \tfrac12\,\mathbb{E}\!\left[(X-\mu)'\Sigma^{-1}(X-\mu)\right].
\]
La dernière espérance se calcule par l'astuce de la trace : un scalaire égale sa
trace, et la trace est cyclique, d'où
\[
\mathbb{E}\!\left[(X-\mu)'\Sigma^{-1}(X-\mu)\right]
= \mathbb{E}\!\left[\operatorname{tr}\!\left(\Sigma^{-1}(X-\mu)(X-\mu)'\right)\right]
= \operatorname{tr}(\Sigma^{-1}\Sigma) = \operatorname{tr}(I_N) = N .
\]
En regroupant, $h(q)=\tfrac12\left[N\log2\pi+\log\det\Sigma+N\right]
=\tfrac12\log\!\left((2\pi e)^N\det\Sigma\right)$.
\end{proof}

\begin{proof}[Étape B : factorisation du déterminant]
Avec $\Sigma=DRD$, $D=\operatorname{diag}(\sigma_1,\dots,\sigma_N)$ et $R$ la
matrice de corrélation, la multiplicativité du déterminant donne
$\det\Sigma=\det(D)\det(R)\det(D)=\left(\prod_i\sigma_i\right)^2\det R$.
\end{proof}

\begin{proof}[Étape C : substitution]
$\tfrac12\log\det\Sigma
=\tfrac12\log\!\left[\left(\prod_i\sigma_i\right)^2\det R\right]
=\sum_i\log\sigma_i+\tfrac12\log\det R$, d'où
$h(q)=\tfrac N2\log(2\pi e)+\sum_i\log\sigma_i+\tfrac12\log\det R$, qui est
l'équation~\eqref{eq:ceiling}.
\end{proof}

\paragraph{Signe du canal de dépendance.} L'inégalité de Hadamard donne
$\det R\le\prod_i R_{ii}=1$ pour une matrice semi-définie positive. Une
démonstration plus parlante passe par le spectre : $\det R=\prod_k\lambda_k$
avec $\sum_k\lambda_k=\operatorname{tr}R=N$, et l'inégalité arithmético-géométrique
donne $\prod_k\lambda_k\le\left(\sum_k\lambda_k/N\right)^N=1$, avec égalité si
et seulement si toutes les valeurs propres valent~1, c'est-à-dire $R=I$. Donc
$H^{dep}\le0$ : \emph{la corrélation détruit toujours de l'entropie}.

Dans le cas bivarié, $\det R=1-\rho^2$ et $H^{dep}=\tfrac12\log(1-\rho^2)$ ; pour
$\rho=0{,}9$ cela vaut $-0{,}83$~nat. Remarquons que $-H^{dep}$ est alors
exactement l'information mutuelle entre les deux actifs : la corrélation est de
l'information partagée, donc de l'entropie détruite.

\paragraph{Note de lecture sur les tableaux.} Les colonnes $H^{cov}$ des
tableaux de résultats de la version arXiv~v1 rapportent $H^{vol}+H^{dep}$ et
\textbf{omettent la constante} $\tfrac N2\log(2\pi e)$, qui vaut $28{,}38$~nats
pour $N=20$. Ce décalage constant est sans conséquence puisque seules les
différences entre régimes et les variations temporelles sont interprétées, mais
il explique que les niveaux rapportés soient négatifs.

\section{Le déficit d'entropie : pourquoi $\DD=h(q)-h(p)$}
\label{app:deficit}

La divergence de Kullback--Leibler entre une loi $p$ et une loi de référence $q$
est $\KL{p}{q}=\int p\log(p/q)$, positive et nulle si et seulement si $p=q$
presque partout. Elle n'est ni symétrique ni munie d'une inégalité triangulaire :
ce n'est pas une distance, d'où le terme de divergence. Les lettres \og KL \fg{}
sont les initiales de Kullback et Leibler.

\begin{proposition}
\label{prop:deficit}
Si $q=\mathcal N(\mu,\Sigma)$ a \emph{les deux premiers moments de $p$}, alors
$\KL{p}{q}=h(q)-h(p)$.
\end{proposition}

\begin{proof}
$\KL{p}{q}=-h(p)-\mathbb{E}_p[\log q(X)]$. Or $\log q$ est exactement
quadratique en $x$, donc son espérance ne dépend que de la moyenne et de la
covariance de la loi sous laquelle on l'évalue. Comme $p$ et $q$ partagent ces
deux moments,
$\mathbb{E}_p\!\left[(X-\mu)'\Sigma^{-1}(X-\mu)\right]=\operatorname{tr}(\Sigma^{-1}\Sigma)=N
=\mathbb{E}_q[\cdot]$, d'où $\mathbb{E}_p[\log q]=\mathbb{E}_q[\log q]=-h(q)$ et
donc $\KL{p}{q}=h(q)-h(p)$.
\end{proof}

L'hypothèse de coïncidence des moments est essentielle : sans elle l'identité
tombe. C'est tout le contenu de la proposition~\ref{prop:maxent}, lue comme une
mesure plutôt que comme un théorème d'existence. Économiquement, $\DD$ est le
prix en nats du mensonge \og le marché est gaussien \fg{} : l'incertitude
apparente que la description gaussienne surestime, parce que la loi véritable est
plus concentrée que ne le laisse croire sa seule covariance.

\section{L'indice structurel est une divergence}
\label{app:index}

\begin{proposition}
\label{prop:index}
Avec $q_i=\mathcal N(\mu_i,\sigma_i^2)$ les marginales gaussiennes de mêmes
échelles que $p$,
\[
\KL{p}{\textstyle\prod_i q_i} \;=\; \DD - H^{dep} \;=\; \JJ .
\]
\end{proposition}

\begin{proof}
Par définition,
\[
\begin{aligned}
\KL{p}{\textstyle\prod_i q_i} &= -h(p) - \mathbb{E}_p\!\left[\textstyle\sum_i\log q_i(X_i)\right]
= -h(p) + \sum_i\left[\tfrac12\log(2\pi\sigma_i^2)+\tfrac12\right]\\
&= -h(p) + \sum_i\log\sigma_i + \tfrac N2\log(2\pi e),
\end{aligned}
\]
où l'on a utilisé $\mathbb{E}_p[(X_i-\mu_i)^2]=\sigma_i^2$, vrai par
construction. En y injectant l'équation~\eqref{eq:three}, soit
$h(p)=\tfrac N2\log(2\pi e)+\sum_i\log\sigma_i+H^{dep}-\DD$, les termes
d'échelle et la constante se télescopent et il reste $\DD-H^{dep}$.
\end{proof}

C'est une identité exacte, non une approximation, et c'est elle qui justifie la
définition~\ref{def:J}. Retirer le canal d'échelle de \eqref{eq:three} pourrait
passer pour un bricolage ; l'objet qui reste se trouve être une divergence de
Kullback--Leibler en bonne et due forme, et hérite donc gratuitement de trois
propriétés : elle est positive ; elle est nulle si et seulement si les rendements
sont indépendants \emph{et} gaussiens, ce qui fournit un point de référence bien
défini plutôt qu'arbitraire ; et elle est invariante par changement d'unités
actif par actif, puisque $p$ et $\prod_i q_i$ se transforment de la même façon.

Il est utile de rapprocher cette identité de la règle de chaînage générale
$\KL{p}{\prod_i q_i}=\sum_i\KL{p_i}{q_i}+I(x)$, où
$I(x)=\KL{p}{\prod_i p_i}$ est la multi-information. Le premier terme est la
somme des déficits marginaux $\sum_i\Dun(x_i)$, le second la dépendance totale.
L'estimateur employé dans l'article estime le premier terme par l'approximation
de \citet{hyvarinen1998} et remplace le second par sa part linéaire
$-H^{dep}$ ; la dépendance de queue non linéaire est donc omise de $\JJ$ et
mesurée séparément par un taux de co-dépassement.

\section{Information mutuelle : les deux propositions et un exemple}
\label{app:mi}

\paragraph{Démonstration de la proposition~\ref{prop:conditioning}.} Par
définition de l'entropie conditionnelle,
$\mathbb{E}_Y[h(X\mid Y)]=h(X,Y)-h(Y)$, et par
\eqref{eq:mi}, $I(X;Y)=h(X)+h(Y)-h(X,Y)$. En sommant les deux,
$\mathbb{E}_Y[h(X\mid Y)]=h(X)-I(X;Y)$. La positivité de $I$ découle de celle de
la divergence de Kullback--Leibler dont elle est un cas particulier, avec
égalité si et seulement si $p_{XY}=p_X\otimes p_Y$.

\paragraph{Démonstration de la proposition~\ref{prop:constraint}.} Immédiate :
une contrainte supplémentaire ne peut que réduire l'ensemble réalisable, et le
supremum d'une fonctionnelle sur un ensemble plus petit ne peut pas augmenter.

\paragraph{Pourquoi les deux ne sont pas équivalents.} La première quantifie
mais ne vaut qu'en moyenne ; la seconde vaut réalisation par réalisation mais
porte sur le plafond et ne quantifie rien. Un contre-exemple discret suffit à
montrer qu'une réalisation particulière peut \emph{augmenter} l'entropie. Soit
$X$ binaire avec $\mathbb{P}(X=a)=0{,}99$, d'entropie $H(X)=0{,}056$~nat. Soit
un signal $Y$ qui, avec probabilité $0{,}98$, lève complètement le doute
($H(X\mid Y=0)=0$), et avec probabilité $0{,}02$ conduit à
$\mathbb{P}(X=a\mid Y=1)=1/2$, soit $H(X\mid Y=1)=\log 2=0{,}693$~nat --- douze
fois l'entropie a priori. La moyenne vaut néanmoins
$0{,}98\times0+0{,}02\times0{,}693=0{,}014$~nat, bien en dessous de $0{,}056$,
et $I(X;Y)=0{,}042$~nat. C'est exactement la raison pour laquelle l'article
prédit une composante de saut \emph{principalement}, et non exclusivement,
unilatérale.

\paragraph{Exemple gaussien chiffré.} Soit $X\sim\mathcal N(0,\sigma^2)$ le
rendement de demain et $Y=X+\varepsilon$ un signal bruité,
$\varepsilon\sim\mathcal N(0,\sigma_\varepsilon^2)$ indépendant. Alors
\begin{equation}
\label{eq:mi-gauss}
I(X;Y)=\tfrac12\log\!\left(1+\frac{\sigma^2}{\sigma_\varepsilon^2}\right),
\qquad
\operatorname{Var}(X\mid Y)=\frac{\sigma^2\sigma_\varepsilon^2}{\sigma^2+\sigma_\varepsilon^2},
\end{equation}
la variance a posteriori ne dépendant pas de la réalisation $y$. Avec
$\sigma=1\%$ par jour :

\begin{center}
\small
\begin{tabular}{@{}lrrrr@{}}
\toprule
Qualité du signal & $\sigma_\varepsilon$ & $I(X;Y)$ (nats) & vol.\ a posteriori & équivalent $\sqrt{2\eta}$ \\
\midrule
tranchant   & $0{,}5\%$ & $0{,}805$ & $0{,}45\%$ & $1{,}27\sigma$ \\
moyen       & $1{,}0\%$ & $0{,}347$ & $0{,}71\%$ & $0{,}83\sigma$ \\
vague       & $3{,}0\%$ & $0{,}053$ & $0{,}95\%$ & $0{,}32\sigma$ \\
\bottomrule
\end{tabular}
\end{center}

Un facteur quinze sépare la taille du saut d'entropie provoqué par une nouvelle
qui tranche de celui provoqué par une rumeur. La dernière colonne traduit chaque
saut dans l'échelle de sévérité de la section~\ref{sec:stress} : c'est le
mouvement directionnel qu'un scénario de même rayon produirait.

\renewcommand{\refname}{Références}
\begin{thebibliography}{99}
\small

\bibitem[Ahmadi-Javid(2012)]{ahmadijavid2012} Ahmadi-Javid, A. (2012). Entropic Value-at-Risk: A New Coherent Risk Measure. \textit{Journal of Optimization Theory and Applications}, 155(3), 1105--1123.

\bibitem[Almog \& Shmueli(2019)]{almog2019} Almog, A., \& Shmueli, E. (2019). Structural Entropy: Monitoring Correlation-Based Networks Over Time With Application To Financial Markets. \textit{Scientific Reports}, 9, 10832.

\bibitem[Ané \& Geman(2000)]{ane2000} Ané, T., \& Geman, H. (2000). Order Flow, Transaction Clock, and Normality of Asset Returns. \textit{Journal of Finance}, 55(5), 2259--2284.

\bibitem[Artzner et al.(1999)]{artzner1999} Artzner, P., Delbaen, F., Eber, J.-M., \& Heath, D. (1999). Coherent Measures of Risk. \textit{Mathematical Finance}, 9(3), 203--228.

\bibitem[Barndorff-Nielsen \& Shephard(2001)]{bns2001} Barndorff-Nielsen, O. E., \& Shephard, N. (2001). Non-Gaussian Ornstein--Uhlenbeck-based Models and Some of Their Uses in Financial Economics. \textit{Journal of the Royal Statistical Society B}, 63(2), 167--241.

\bibitem[Barndorff-Nielsen \& Shephard(2004)]{bns2004} Barndorff-Nielsen, O. E., \& Shephard, N. (2004). Power and Bipower Variation with Stochastic Volatility and Jumps. \textit{Journal of Financial Econometrics}, 2(1), 1--37.

\bibitem[Black(1976)]{black1976} Black, F. (1976). Studies of Stock Price Volatility Changes. \textit{Proceedings of the Business and Economic Statistics Section, American Statistical Association}, 177--181.

\bibitem[Christie(1982)]{christie1982} Christie, A. A. (1982). The Stochastic Behavior of Common Stock Variances: Value, Leverage and Interest Rate Effects. \textit{Journal of Financial Economics}, 10(4), 407--432.

\bibitem[Doob(1942)]{doob1942} Doob, J. L. (1942). The Brownian Movement and Stochastic Equations. \textit{Annals of Mathematics}, 43(2), 351--369.

\bibitem[Engle \& Ng(1993)]{engle1993} Engle, R. F., \& Ng, V. K. (1993). Measuring and Testing the Impact of News on Volatility. \textit{Journal of Finance}, 48(5), 1749--1778.

\bibitem[Nelson(1991)]{nelson1991} Nelson, D. B. (1991). Conditional Heteroskedasticity in Asset Returns: A New Approach. \textit{Econometrica}, 59(2), 347--370.

\bibitem[Bollerslev \& Ghysels(1996)]{bollerslev1996} Bollerslev, T., \& Ghysels, E. (1996). Periodic Autoregressive Conditional Heteroscedasticity. \textit{Journal of Business \& Economic Statistics}, 14(2), 139--151.

\bibitem[Breuer \& Csiszár(2013)]{breuer2013} Breuer, T., \& Csiszár, I. (2013). Systematic Stress Tests with Entropic Plausibility Constraints. \textit{Journal of Banking \& Finance}, 37(5), 1552--1559.

\bibitem[Caraiani(2018)]{caraiani2018} Caraiani, P. (2018). Modeling the Comovement of Entropy between Financial Markets. \textit{Entropy}, 20(6), 417.

\bibitem[Clark(1973)]{clark1973} Clark, P. K. (1973). A Subordinated Stochastic Process Model with Finite Variance for Speculative Prices. \textit{Econometrica}, 41(1), 135--155.

\bibitem[Cover \& Thomas(2006)]{cover2006} Cover, T. M., \& Thomas, J. A. (2006). \textit{Elements of Information Theory}, 2\ieme{} éd. Wiley.

\bibitem[Csiszár(1975)]{csiszar1975} Csiszár, I. (1975). $I$-Divergence Geometry of Probability Distributions and Minimization Problems. \textit{Annals of Probability}, 3(1), 146--158.

\bibitem[Diebold \& Mariano(1995)]{dm1995} Diebold, F. X., \& Mariano, R. S. (1995). Comparing Predictive Accuracy. \textit{Journal of Business \& Economic Statistics}, 13(3), 253--263.

\bibitem[Engle(2002)]{engle2002} Engle, R. (2002). Dynamic Conditional Correlation. \textit{Journal of Business \& Economic Statistics}, 20(3), 339--350.

\bibitem[Forbes \& Rigobon(2002)]{forbes2002} Forbes, K. J., \& Rigobon, R. (2002). No Contagion, Only Interdependence: Measuring Stock Market Comovements. \textit{Journal of Finance}, 57(5), 2223--2261.

\bibitem[Holló, Kremer \& Lo Duca(2012)]{hollo2012} Holló, D., Kremer, M., \& Lo Duca, M. (2012). CISS -- A Composite Indicator of Systemic Stress in the Financial System. \textit{ECB Working Paper} No.~1426.

\bibitem[Kritzman et al.(2011)]{kritzman2011} Kritzman, M., Li, Y., Page, S., \& Rigobon, R. (2011). Principal Components as a Measure of Systemic Risk. \textit{Journal of Portfolio Management}, 37(4), 112--126.

\bibitem[Glasserman \& Xu(2014)]{glasserman2014} Glasserman, P., \& Xu, X. (2014). Robust Risk Measurement and Model Risk. \textit{Quantitative Finance}, 14(1), 29--58.

\bibitem[Hyvärinen(1998)]{hyvarinen1998} Hyvärinen, A. (1998). New Approximations of Differential Entropy for Independent Component Analysis and Projection Pursuit. \textit{Advances in Neural Information Processing Systems}, 10, 273--279. (Constantes telles que dans Hyvärinen \& Oja, 2000, \textit{Neural Networks} 13, 411--430.)

\bibitem[Jaynes(1957)]{jaynes1957} Jaynes, E. T. (1957). Information Theory and Statistical Mechanics. \textit{Physical Review}, 106, 620--630.

\bibitem[Kyle(1985)]{kyle1985} Kyle, A. S. (1985). Continuous Auctions and Insider Trading. \textit{Econometrica}, 53(6), 1315--1335.

\bibitem[Ledoit \& Wolf(2004)]{ledoit2004} Ledoit, O., \& Wolf, M. (2004). A Well-Conditioned Estimator for Large-Dimensional Covariance Matrices. \textit{Journal of Multivariate Analysis}, 88(2), 365--411.

\bibitem[Mandelbrot \& Taylor(1967)]{mandelbrot1967} Mandelbrot, B., \& Taylor, H. M. (1967). On the Distribution of Stock Price Differences. \textit{Operations Research}, 15(6), 1057--1062.

\bibitem[Monroe(1978)]{monroe1978} Monroe, I. (1978). Processes That Can Be Embedded in Brownian Motion. \textit{Annals of Probability}, 6(1), 42--56.

\bibitem[Papla \& Siedlecki(2024)]{papla2024} Papla, D., \& Siedlecki, R. (2024). Entropy as a Tool for the Analysis of Stock Market Efficiency During Periods of Crisis. \textit{Entropy}, 26(12), 1079.

\bibitem[Politis \& Romano(1994)]{politis1994} Politis, D. N., \& Romano, J. P. (1994). The Stationary Bootstrap. \textit{Journal of the American Statistical Association}, 89(428), 1303--1313.

\bibitem[Tauchen \& Pitts(1983)]{tauchen1983} Tauchen, G. E., \& Pitts, M. (1983). The Price Variability--Volume Relationship on Speculative Markets. \textit{Econometrica}, 51(2), 485--505.

\bibitem[Vasicek(1976)]{vasicek1976} Vasicek, O. (1976). A Test for Normality Based on Sample Entropy. \textit{Journal of the Royal Statistical Society B}, 38(1), 54--59.

\bibitem[Zhou, Cai \& Tong(2013)]{zhou2013} Zhou, R., Cai, R., \& Tong, G. (2013). Applications of Entropy in Finance: A Review. \textit{Entropy}, 15(11), 4909--4931.

\end{thebibliography}

\end{document}
